\section{Architecture}\label{sec:arch}

\subsection{Four-tier ANU fallback}\label{sec:arch-anu}

The \qmirror{} entropy stack defines four fallback tiers, evaluated
in order. Figure~\ref{fig:arch} sketches the data flow.

\begin{figure}[t]
\centering
% PLACEHOLDER: replace with TikZ block diagram or PDF export
\begin{lstlisting}[basicstyle=\ttfamily\tiny,frame=single,framesep=4pt]
*** Figure 1 placeholder (text-flowchart) ***

+----------------------------+
| qmirror.api               |
| (qrng/chsh/iit/circuit)   |
+-------------+--------------+
              |
   amplitudes p_i
              v
+----------------------------+         +-----------------+
| qmirror.engine (Aer|Cirq) |  <----- | exact unitary U |
+-------------+--------------+         +-----------------+
              |
       wants randomness
              v
+----------------------------+
| qmirror.entropy (4-tier)  |
|  T0 HMAC-DRBG (NIST 90A)  |  <----- ANU seed (boot)
|  T1 ANU keyed REST        |  <----> https x-api-key
|  T2 ANU keyless legacy    |  <----> https courtesy
|  T3 mock LCG              |  <----- NEXUS_QMIRROR_MOCK=1
+-------------+--------------+
              |
            bytes
              v
+----------------------------+
| qmirror.sampler           |
| (CDF inversion)           |
+-------------+--------------+
              |
       outcome index
              v
       (back to api)
\end{lstlisting}
\caption{Architecture of the \qmirror{} substrate. Tier~0 (HMAC-DRBG
SHA-256, NIST SP~800-90A keyed by ANU bytes) is the deterministic
regression default; Tier~1 (ANU QRNG keyed REST API) is the live
quantum-entropy path; Tier~2 (ANU QRNG keyless legacy endpoint) is the
courtesy fallback; Tier~3 (mock LCG) is the CI-deterministic
invariance path. Aer (or Cirq) provides exact unitary evolution; the
entropy stream feeds CDF inversion at the measurement collapse step.}
\label{fig:arch}
\end{figure}

\begin{lstlisting}[caption={Tier definitions (extracted from
nexus/.roadmap.qmirror)},label={lst:tiers}]
T0: HMAC-DRBG SHA-256 keyed by ANU bytes (cached seed, 32 KB ring)
    - Default in CI; deterministic given seed; NIST SP 800-90A
T1: ANU QRNG REST API (live, key-gated, ~10^6 byte/req)
    - api.quantumnumbers.anu.edu.au with x-api-key header
T2: ANU QRNG REST API (legacy keyless endpoint)
    - qrng.anu.edu.au; courtesy ~1 req/min; backup if T1 quota exhausted
T3: Mock LCG (deterministic linear congruential)
    - Activated by NEXUS_QMIRROR_MOCK = 1; cross-module byte-identity
      with nexus/modules/qrng/mock_qrng.hexa for regression invariance
\end{lstlisting}

Tier promotion is automatic on HTTP failure or rate-limit response;
tier demotion requires explicit env-var override. This pattern was
hardened by \texttt{nexus@a962c4c81} (4-tier ANU revision) and
\texttt{nexus@02225e87} (JSON whitespace tolerance fix that enabled
F1 LIVE in-band PASS).

\subsection{HMAC-DRBG SHA-256}

The deterministic regression path uses HMAC-DRBG (SHA-256) per NIST
SP~800-90A Rev.~1~\cite{NIST-SP-800-90A}, seeded by an ANU QRNG fetch
at session boot. This provides:

\begin{itemize}
\item Determinism for CI regression: a fixed seed reproduces a fixed
  bit stream.
\item Quantum-derived entropy at session origin: the seed is genuinely
  quantum.
\item NIST SP~800-22 statistical strength on the expanded stream
  (validated by cond.4).
\end{itemize}

The reseeding cadence is operator-configurable; the default is
once-per-process. Quarterly IonQ refresh as anchor is documented in
the spec section~10 M1.

\subsection{Qiskit Aer engine bridge}

Aer is Python-only; the \texttt{nexus} codebase is hexa-strict (raw\#9:
zero \texttt{.py} files on Mac side). \qmirror{} resolves this via a
documented \textbf{\texttt{python\_bridge} concession}: a single
subdirectory \texttt{nexus/modules/qmirror/\_python\_bridge/} contains
all and only the \texttt{.py} shims required to invoke Aer
(\texttt{aer\_runner.py}), \texttt{pyphi}
(\texttt{iit\_mip\_runner.py}), and the Braket helper. The hexa-side
modules spawn these as subprocesses via JSON over stdin/stdout. There
is no in-process linking; the FSF Mere Aggregation
doctrine~\cite{FSF-MereAggregation} applies (relevant for license
isolation; see Section~\ref{sec:limitations}).

\subsection{IIT 4.0 MIP via pyphi}

\texttt{qmirror.iit\_mip.calc(tpm, partition\_hint)} accepts a
state-by-node TPM in \texttt{pyphi}'s expected shape and invokes
\texttt{pyphi.sia()}. Output schema matches the reference
\texttt{braket\_iit40\_mip\_2026\_05\_02 verdict.json}:
\texttt{phi\_star} (\texttt{float}), \texttt{partition\_used}
(\texttt{list[int]}), \texttt{mip\_sec} (\texttt{float}), \texttt{ok}
(\texttt{int}).

For $n \le 6$ the search is exact; for $6 < n \le 12$ the
\texttt{CUT\_ONE\_APPROXIMATION} heuristic is used; for $n > 12$ the
call raises \texttt{E\_PHI\_TOO\_LARGE}. The \texttt{pyphi} version
pin (4.0 \texttt{feature/iit-4.0} branch commit \texttt{b78d0e3}) is
\textbf{load-bearing}: newer \texttt{pyphi} may change MIP search
heuristics and produce drift that looks like substrate change but is
software version drift.

\subsection{CHSH circuit (canonical geometry)}

\texttt{qmirror.chsh.run(n\_trials)} builds four Bell-state circuits
matching the schema in
\texttt{nexus\_chsh\_bell\_2026\_05\_02/circuit\_*.json}. The canonical
geometry uses an $R_{y}(-\theta)$ rotation with the four standard
angles $(0, \pi/2, \pi/4, -\pi/4)$, yielding the textbook
$S = 2\sqrt{2} \approx 2.828$ quantum bound. \qmirror{} executes each
circuit on Aer state-vector, draws $n_{\text{trials}}$ measurement
outcomes per circuit using ANU-fed CDF inversion, computes the four
correlators, and returns $S$, $\sigma_{S}$, and the per-correlator
$E$ values. Expected output band:
$S \in [2.7, 2.85]$ at $n_{\text{trials}} = 1000$.

\subsection{Module surface}

\begin{lstlisting}[caption={\qmirror{} module layout},label={lst:modules}]
nexus/modules/qmirror/
  entropy.hexa       T1/T2/T3 ANU REST + cache + mock
  sampler.hexa       amplitude -> outcome via entropy stream
  engine_aer.hexa    Aer subprocess shim
  engine_cirq.hexa   Cirq alt (scaffolded)
  circuit.hexa       circuit.exec() public API
  qrng.hexa          drop-in replacement for HMAC-DRBG
  chsh.hexa          Bell test (S, sigma, violation)
  iit_mip.hexa       pyphi sia() shim -> phi_star, MIP partition
  tomography.hexa    process tomography (qmirror 2.0 axis)
  phi.hexa           anima_phi_v3_canonical port (qmirror 2.0)
  selftest.hexa      __QMIRROR__ {PASS|FAIL}
  _python_bridge/    aer_runner.py, iit_mip_runner.py, ...
\end{lstlisting}
